\pdfoutput=1
\documentclass[11pt]{article}
\usepackage[T1]{fontenc}
\usepackage{lmodern,microtype}
\setlength{\emergencystretch}{2em}
\usepackage[margin=1in]{geometry}
\usepackage{amsmath,amssymb,amsthm,mathtools}
\usepackage[hidelinks]{hyperref}
\hypersetup{pdftitle={What a Retained-Mode Observer Can Recover in Quantum Einstein--Rosen Dynamics}}
\newtheorem{theorem}{Theorem}
\newtheorem{proposition}[theorem]{Proposition}
\newcommand{\cH}{\mathcal H}
\newcommand{\cT}{\mathcal T}
\newcommand{\cB}{\mathcal B}
\newcommand{\cI}{\mathcal I}
\newcommand{\cU}{\mathcal U}
\newcommand{\tr}{\operatorname{Tr}}
\newcommand{\id}{\mathrm{id}}
\newcommand{\dd}{\mathrm d}
\newcommand{\one}{\mathbf 1}
\newcommand{\vac}{v_0}
\newcommand{\ketbra}[2]{|#1\rangle\langle #2|}
\title{What a Retained-Mode Observer Can Recover\\ in Quantum Einstein--Rosen Dynamics}
\author{}
\date{}
\begin{document}
\maketitle

\begin{abstract}
In the exact polarized Einstein--Rosen quantization, low radial modes evolve in asymptotic physical time with a phase controlled by the conserved auxiliary energy of discarded high modes. We identify the complete information about a normal joint preparation that is visible to low-mode instruments: a positive trace-class-valued measure on the spectrum of the exponential of the high auxiliary Hamiltonian. Equality of this measure is equivalent to equality of every finite retained history. Necessity uses just an initial effect and reset, one vacuum--packet Ramsey gap, and exact elapsed-time data in an arbitrarily small interval. We use bounded full-space effects, so no sharp continuum energy eigenstate or postselection on a packet span is needed. An explicit affine entangled preparation, a positive two-reset memory witness, a fixed-cut obstruction to finite-dimensional memory, and equal-mean high-energy laws that cannot be related by a deterministic change of clock illustrate the result. The classification is exact on the physical Fock space, not a perturbative expansion or a claim of spatial factorization for gravity.
\end{abstract}

\section{Physical Hamiltonian and the operational question}

Consider the vacuum polarized cylindrical Einstein--Rosen model in the asymptotic-time quantization of Ashtekar and Pierri \cite[\S4.2, Eq.~(37)]{AP}. Its physical Hilbert space is the symmetric Fock space
\[
  \cH=\Gamma_{\rm s}\bigl(L^2(\mathbb R_+,\dd k)\bigr),\qquad
  H_0=\dd\Gamma(k),\qquad
  H=A(\one-e^{-bH_0}),\qquad A,b>0 .
\]
We set \(\hbar=1\) and keep the positive constants \(A,b\) to allow for the normalization of the radial measure and coupling used in the different conventions. The operator \(H_0\) is vacuum-normal-ordered. The time \(t\) throughout is generated by \(H\) and normalized at infinity, not the free auxiliary time generated by \(H_0\). We do not choose a separate clock for each input state.

Fix \(\Lambda>0\) and use the one-particle spectral split \((0,\Lambda)\sqcup(\Lambda,\infty)\), up to a null set. Fock factorization over this orthogonal direct sum and spectral calculus give
\begin{gather}
  \cH=\cH_L\otimes\cH_H,\qquad
  H_0=H_{0,L}\otimes\one+\one\otimes H_{0,H},\nonumber\\
  q=e^{-bH_{0,L}},\qquad R=e^{-bH_{0,H}},\qquad
  H=A(\one-q\otimes R).
  \label{eq:physical}
\end{gather}
The exponential \(R\) is a bounded positive contraction and has no eigenprojection at zero. This particular split is by reduced \emph{radial modes}, not by spatial regions or subregions of an unfixed gravitational phase space. There is a bona fide Fock tensor product for this split. A retained instrument has normal completely positive outcome maps on the low factor, trace nonincreasing, whose sum is trace preserving. They may act coherently across retained Fock particle numbers; this is an allowed control in the question, not an inferred capability of point-local detectors. Only those maps and physical-time evolution may be used after initialization. In particular a high-mode measurement or an intervention that changes its Hamiltonian is not in the class being classified.

For a trace-one positive trace-class operator \(\rho\) on \(\cH\), write \(P_B=E_R(B)\) for the spectral projections of \(R\) and set
\begin{equation}
 M_\rho(B)=\tr_H\!\left[(\one\otimes P_B)\rho(\one\otimes P_B)\right],
 \qquad B\subset[0,1]\ \hbox{Borel}.
 \label{eq:measure}
\end{equation}
Its total mass is the reduced low density \(M_\rho([0,1])=\rho_L\). The problem is whether this whole measure, a smaller record, or the full joint state is needed to predict the probabilities of all retained interventions.

\begin{theorem}[Operational classification]\label{thm:main}
For any two normal joint density operators on the physical Fock space, the following are equivalent.
\begin{enumerate}
\item Their measures \(M_\rho\) in \eqref{eq:measure} agree on every Borel set.
\item Every finite retained instrument history in physical time has the same probabilities.
\item For every retained effect used in an initial effect-and-reset operation, the two preparations agree on one fixed vacuum--packet Ramsey protocol for times \(0<t<\epsilon\), where \(\epsilon>0\) may be arbitrarily small. The packet is any fixed normalized smooth one-particle mode supported in a compact subinterval of \((0,\Lambda)\), and both reading quadratures may be used.
\end{enumerate}
The history probabilities are realized by a positive, conserved spectral parameter and no reset of the discarded factor. The classification is an exact statement about operational equivalence, not a finite-data reconstruction guarantee.
\end{theorem}

Barbero, Garay and Villase\~nor \cite[\S\S4--6]{BGV} already use the bounded nonlinear Hamiltonian and compute the exact multiparticle energy phases and correlation functions; the phase in \eqref{eq:physical} is not new. In the broader setting of quantum multitime processes \cite{PollockProcess,PollockMarkov}, Goswami et al.\ \cite[Thm.~1]{Goswami} derive classical common-cause instrument histories from a conserved discrete environment basis, including initially correlated states. We use the corresponding exact spectral mechanism on the continuous, untruncated gravitational model and prove the converse for an explicit retained experiment. This identifies which initial-state data are necessary, not just one representation that suffices.

\section{The positive history law}\label{sec:law}

\subsection{Trace-class measure and fiber formula}

Let \(\cT(\cH_L)\) denote the trace-class operators and \(\cB(\cH_L)\) the bounded operators. For any \(X\in\cB(\cH_L)\), trace cyclicity gives
\begin{equation}
 \tr_L[X M_\rho(B)]=\tr[(X\otimes P_B)\rho].
 \label{eq:pair}
\end{equation}
This justifies an equivalent partial trace with a single projector and does not perform a measurement on the high factor. The trace of \(M_\rho\) is the scalar probability measure \(\mu(B)=\tr[(\one\otimes P_B)\rho]\). Positivity and finite additivity in \eqref{eq:pair}, followed by the fact that the trace norm of a positive remainder equals its trace, prove countable additivity in trace norm. In particular \(M_\rho\) has bounded variation.

There is a trace-one positive conditional density \(D_\rho(r)\) for \(\mu\)-almost every \(r\). For a direct proof, fix a countable orthonormal low basis, take scalar Radon--Nikodym derivatives of each matrix entry with respect to \(\mu\), and impose all finite rational-vector positivity tests on a common full-measure set. The diagonal derivatives are nonnegative and sum to one there by monotone convergence. The resulting positive matrix defines a positive trace-class operator of trace one. Its finite compressions converge in trace norm pointwise and are measurable, hence the field is trace-norm measurable. Integrating in trace norm gives \(M_\rho(\dd r)=D_\rho(r)\mu(\dd r)\). No energy moment of \(\rho\) is assumed.

Define the fixed-\(r\) evolution on low states by
\begin{equation}
 u_r(t)=e^{iAtrq},\qquad \cU_r(t)(S)=u_r(t)S u_r(t)^*.
 \label{eq:fiber}
\end{equation}
The common scalar \(e^{-iAt}\) in the physical propagator drops from conjugations. Let \(t_j-t_{j-1}=\Delta t_j\geq 0\) be successive durations and let \(\cI_{j,o_j}\) be the selected instrument outcomes. Their history law is
\begin{equation}
 \Pr_\rho(o_1,\ldots,o_n)
 =\int_{[0,1]}\!\!
 \tr_L\!\left[
 \cI_{n,o_n}\cU_r(\Delta t_n)\cdots
 \cI_{1,o_1}\cU_r(\Delta t_1)\,M_\rho(\dd r)
 \right].
 \label{eq:history}
\end{equation}
One can use the conditional density above or a bounded-variation integral with norm-continuous operator coefficient. Maps and times can depend on previous outcomes; then along the selected branch the indicated maps and times are fixed.

To prove \eqref{eq:history}, for a norm-continuous bounded-operator field \(X(r)\) let
\[
 \mathfrak D[X]=\int_{[0,1]} X(r)\otimes E_R(\dd r).
\]
Construct it by uniform approximation with simple fields. Orthogonality of the high projectors gives the sup norm bound, and \eqref{eq:pair} gives its expectation as the pairing with \(M_\rho\). The low Heisenberg map of an outcome acts fiberwise on \(\mathfrak D[X]\) by its normal bounded adjoint, first on simple fields and then by continuity. The physical unitary acts by \(X(r)\mapsto u_r(t)^*X(r)u_r(t)\) under the same norm approximation. Iterating these two facts backward from \(X(r)=\one\) proves the formula without a diagonal density of the high state or a continuum eigenvector. A retained spectator initially included in the low register is allowed by the identical proof with \(q\otimes\one_{\rm sp}\). If it is initially independent, its fixed density simply tensors with \(M_\rho\); initially correlated spectators require the corresponding enlarged measure.

The law is positive and normalized because each term at fixed \(r\) is a positive quantum instrument history. Summing final outcomes recovers the shorter history; summing an interior outcome leaves its trace-preserving nonselective operation in place. Subdividing a time evolution without an intervention does not alter anything, since \(\cU_r(s+t)=\cU_r(s)\cU_r(t)\). For each recorded outcome the unnormalized conditional record is propagated by \(\cI_{j,o_j}\cU_r(\Delta t_j)\) on the same \(r\); the normalization for the conditioned record is the integrated outcome probability. Thus the hidden parameter does not change with subdivision or reset. Its posterior distribution can change. Formula \eqref{eq:history} also proves the first implication of Theorem~\ref{thm:main}.

\subsection{Multiple discarded bands and cutoff removal}

For two disjoint discarded radial bands with high exponential operators \(R_1,R_2\), form \eqref{eq:measure} using their joint spectral projections with the final retained low factor. Their physical total control is \(R_1R_2\). Spectral calculus and \eqref{eq:pair} give the combined measure as the pushforward of the joint measure under \((r_1,r_2)\mapsto r_1r_2\). This statement also follows from successive elimination if the intermediate retained register includes the not-yet-discarded second band: taking its spectral projection and trace in the resulting operator measure gives the same rectangle measures. Finite additivity and then uniqueness as measures give the identity for all Borel sets. The same argument applies to a finite number of bands. These are equalities for histories acting only on the finally retained factor, not a statement of equivalence after real interventions on an eventually discarded band. Replacing the joint record by a product of its marginals is not this pushforward.

For clarity, the formulas do not require a particle cutoff. As one useful compatible removal procedure, truncate the one-particle multiplier \(k\) to \(k\,\one_{k\leq N}\) while leaving the Hilbert space fixed. On each particle sector the modified auxiliary energies converge pointwise. Their bounded exponentials converge strongly by dominated convergence, also on the full direct sum. The bounded physical propagators converge strongly with adjoints, uniformly on compact time intervals: on each fixed spectral tuple use \(\sup_{|t|\leq T}|e^{-itx_N}-e^{-itx}|\leq \min(2,T|x_N-x|)\). Conjugation on a trace-class operator then converges in trace norm by reduction to finite rank. One can also replace \(\rho\) by \(P_N\rho P_N/\tr(P_N\rho)\) for projections \(P_N\uparrow\one\), which converge to \(\rho\) in trace norm. Normal completely positive trace-nonincreasing maps contract the trace norm on self-adjoint trace-class input. A finite telescoping bound therefore gives convergence of each probability for fixed instruments and a fixed finite history, with the same initial state or such projecting approximants. This is not a uniform-in-history estimator or a claim about arbitrary regulator schemes.

\section{Recovery by a physical vacuum--packet gap}\label{sec:recovery}

Let \(F\) be any retained effect, \(0\leq F\leq \one\), and denote the finite positive scalar measure by \(m_F(\dd r)=\tr_L(FM_\rho(\dd r))\). Take the initial selected outcome map to be
\begin{equation}
 S\longmapsto \tr_L(FS)\,\ketbra{\psi_j}{\psi_j},\qquad
 \psi_j=\frac{\vac+v_j}{\sqrt2}.
 \label{eq:reset}
\end{equation}
Here \(\vac\) is the low Fock vacuum, and \(v_j\) is a normalized smooth one-particle packet with wavefunction \(v_j(k)\) of compact support in \((0,\Lambda)\). The complementary outcome can also reset to \(\psi_j\), with \(\one-F\) as its effect, so this is a genuine instrument. The high input is not reset; the selected low operation reads the unnormalized \(m_F\) in \eqref{eq:history}.

Define \(X_j=\ketbra{\vac}{v_j}+\ketbra{v_j}{\vac}\) and
\(Y_j=i\ketbra{\vac}{v_j}-i\ketbra{v_j}{\vac}\) as operators on the \emph{entire} low Hilbert space, zero on the orthogonal complement of this two-dimensional span. Since these are self-adjoint contractions, the two effects \((\one\pm X_j)/2\) and, in the separate quadrature setting, \((\one\pm Y_j)/2\) form full-space binary readouts. In particular both outcomes assign half weight to population outside the packet span. This retains leakage rather than projecting onto successful survival and renormalizing. Equivalently one can use sharp signed outcomes and retain a separate zero outcome on the complement. In the binary convention the signed outcome differences for the two choices, with the initial selection still unnormalized, give the real and imaginary parts of
\begin{equation}
  G_j(t;F)
   = \int_{[0,1]} m_F(\dd r)
      \int_{0}^\Lambda |v_j(k)|^2 e^{iAt r a(k)}\,\dd k,\qquad
      a(k)=1-e^{-bk}.
 \label{eq:ramsey}
\end{equation}
Indeed \(u_r(t)\vac=e^{iAtr}\vac\), whereas the phase in \(u_r(t)v_j\) is \(e^{iAtr e^{-bk}}\) pointwise in \(k\). The coherence \(\langle \vac|u_r(t)\ketbra{\psi_j}{\psi_j}u_r(t)^*|v_j\rangle\) is one half of the inner integral in \eqref{eq:ramsey}. The chosen signs of \(X_j,Y_j\) give the stated two signed means. No sharp continuum mode or unbounded detector enters.

\begin{proof}[Proof of the converse in Theorem~\ref{thm:main}]
All finite history equality implies equality of the stated protocols. Since \(0\leq r,a(k)\leq1\), differentiation under both finite integrals gives
\begin{equation}
 \partial_t^nG_j(0;F)
   =(iA)^n \alpha_{j,n}\int r^n m_F(\dd r),
 \qquad
 \alpha_{j,n}=\int_0^\Lambda |v_j(k)|^2a(k)^n\,\dd k>0
 \quad(n\geq0).
 \label{eq:moments}
\end{equation}
In particular \(\alpha_{j,0}=1\); a common arbitrarily short positive-time interval fixes every derivative in \eqref{eq:moments} by continuity and differentiability at zero. Equivalently \(G_j\) has an entire continuation, since its frequencies are bounded. A finite measure on \([0,1]\) is determined by its moments: polynomials are uniformly dense in the continuous functions. Hence the packet protocol determines \(m_F\). Varying the effect \(F\) separates the positive trace-class operators \(M_\rho(B)\); rank-one projections already suffice by polarization. The measure therefore agrees for every Borel set. One packet is sufficient; a second smooth packet in a disjoint positive-gap interval gives an independent form factor \(\alpha_{j,n}\) for the same recovered moments. The real quadrature alone is enough for injectivity by the even moments and the homeomorphism \(r\mapsto r^2\) on \([0,1]\). We retain the complex measurement for the clock comparison below.
\end{proof}

The \emph{ideal gap} is an optional limiting description of this physical test. For normalized smooth packets of width \(\delta\) concentrated around \(k_*\in(0,\Lambda)\), the inner integral in \eqref{eq:ramsey} tends to \(e^{iAt r a(k_*)}\) uniformly in \(r\) and in \(t\) on compact sets. Its deviation is bounded by \(|At|\sup_{|k-k_*|\leq\delta}|a(k)-a(k_*)|\). Thus the formal sharp-gap signal is the Fourier transform of \(m_F\) at the known frequency scale \(A a(k_*)t\). The theorem uses a single nonzero-width packet, not this limit. Agreement on a time set with an accumulation point also suffices in the ideal exact-data sense. No bound on statistical conditioning or on the number of finite-precision samples is asserted.

\section{Compatible entangled data and measurable memory}\label{sec:examples}

\subsection{An assignment on a specified low domain}

Choose orthonormal physical low packets \(\ell_0,\ell_1\) and normalized high one-particle packets \(h_0,h_1\) with disjoint compact auxiliary-energy supports inside \((\Lambda,\infty)\). On the low domain
\(\rho_L(p)=p\ketbra{\ell_0}{\ell_0}+(1-p)\ketbra{\ell_1}{\ell_1}\) with \(p\in[1/3,2/3]\), an affine consistent assignment is
\begin{align}
 \rho_p={}&p\ketbra{\ell_0h_0}{\ell_0h_0}
   +(1-p)\ketbra{\ell_1h_1}{\ell_1h_1}\nonumber\\
 &\quad+c\bigl(\ketbra{\ell_0h_0}{\ell_1h_1}
              +\ketbra{\ell_1h_1}{\ell_0h_0}\bigr),
 \qquad 0<c<\sqrt2/3 .
 \label{eq:assignment}
\end{align}
The determinant of its occupied block is \(p(1-p)-c^2>0\). Orthogonality of the high states gives the specified marginal; partial transpose in the high basis has a block with eigenvalue \(-c\) on the crossed pair \(\ell_0h_1,\ell_1h_0\), proving entanglement. This is an assignment on the given convex segment only, not a positive or completely positive assignment on all low states \cite{Pechukas}. It is used once at initialization.

Write \(\mu_j\) for the \(R\)-law of \(h_j\). The measure is
\begin{equation}
 M_{\rho_p}(\dd r)=p\,\mu_0(\dd r)\ketbra{\ell_0}{\ell_0}
 +(1-p)\,\mu_1(\dd r)\ketbra{\ell_1}{\ell_1}.
 \label{eq:assignmentmeasure}
\end{equation}
In fact \(\langle h_1,P_B h_0\rangle=0\) for every \(B\) by the disjoint spectral supports. This shows how a genuinely entangled full preparation can have a classical conserved record in all retained histories. The off-diagonal information about that full state is invisible to the given controls. Within this assignment the two block effects read the two weights and their conditional laws; neither complete high tomography nor the high spectral degeneracies are identified or needed.

\subsection{A break and reset that does not erase the memory}

Condition on either initial block of \eqref{eq:assignmentmeasure}, with a non-point high law \(\mu_j\), and abbreviate it by \(\nu\). Perform two identical Ramsey rounds using one fixed packet \(v\), a fixed small duration \(t>0\), and the \(+\) outcome of \((\one+X)/2\). At the first readout measure and then prepare the same \(\psi=(\vac+v)/\sqrt2\) for \emph{either} outcome. At fixed \(r\) the success probability of one round is
\begin{equation}
 s(r)=\frac{1+\int |v(k)|^2\cos(At r a(k))\,\dd k}{2}.
 \label{eq:success}
\end{equation}
Choose \(At\max_{\operatorname{supp} v}a(k)<\pi\). Then \(s'(r)<0\) for \(r>0\), since all the supported gaps are positive, so \(s\) is strictly decreasing on \([0,1]\). Given \(r\), the fresh low state makes the next success independent of the preceding success. But the two rounds use the \emph{same} \(r\), not independent samples. With \(p_s=\int s\,\dd\nu\),
\begin{equation}
 \Pr(+,+)-p_s^2
   =\int s(r)^2\,\nu(\dd r)
          -\left(\int s(r)\,\nu(\dd r)\right)^2
   =\operatorname{Var}_{\nu}(s)>0.
 \label{eq:excess}
\end{equation}
For the non-point law \(0<p_s<1\); the success probability of the second round differs after \(+\) and \(-\) by \(\operatorname{Var}_\nu(s)/(p_s(1-p_s))\) despite the identical prepared state. This is a failure of the operational memoryless condition \cite{PollockMarkov} in an allowed causal break. It is not inferred from an unnormalized two-point function or a failure of complete positivity. The enlarged process \eqref{eq:history} is positive; resetting only the low state leaves a different posterior over the conserved control.

\subsection{Exact memory across a fixed cut has unbounded rank}

The same fixed interval and reset instrument give an exact memory-size test. For \(2n\) rounds choose a cut just after the reset in round \(n\). For \(0\leq i,j\leq n\), retain successes in the first \(i\) prefix rounds and the first \(j\) suffix rounds, summing over outcomes in all other rounds. These are events in a \emph{single} \(2n\)-round sequence of actual interventions, not experiments of differing duration. Condition on the same initial block as above. The resulting matrix is
\begin{equation}
 P^{(n)}_{ij}=\int s(r)^{\,i+j}\,\nu(\dd r)
 \quad(0\leq i,j\leq n).
 \label{eq:hankel}
\end{equation}
If \(\nu\) has infinite support, as for the nonzero smooth high packet in \eqref{eq:assignment}, the image law of \(s\) also has infinite support by strict monotonicity. The moment matrix \eqref{eq:hankel} is strictly positive definite: its quadratic form is the integral of the modulus-square of a polynomial, which can vanish only for the zero polynomial. Thus its rank is \(n+1\).

Suppose an alternative simulator uses the same retained reset and a memory system of dimension \(d\) at the cut. All information carried across the cut other than the common fresh low state must be represented on that memory; no uncounted outcome record is allowed to affect the future. Each prefix event then gives a subnormalized memory state \(\sigma_i\) and each suffix event gives a bounded effect \(E_j\) on that memory. Since the instrument at the cut is measure and prepare, the fresh low state is a fixed product with \(\sigma_i\), even after postselection. It follows that \(P^{(n)}_{ij}=\tr(E_j\sigma_i)\) and that \(\operatorname{rank}P^{(n)}\leq d^2\). This is the usual linear rank bound for finite-memory realizations \cite{Fanizza}, applied here at one fixed cut so that the simulator need not be stationary. A \(D\)-state classical memory gives the bound \(D\) instead. No uniformly finite memory suffices for all histories of this single physical preparation: a simulator for \(2n\) such rounds needs \(d^2\geq n+1\) at the chosen cut. This says nothing about approximate simulation at fixed accuracy. For \(n=1\) the first nontrivial determinant is already the two-reset variance in \eqref{eq:excess}.

\paragraph{Which physical preparations are memoryless?}
Taking \(F=\one\), the initial reset is deterministic and the law \(\nu\) in \eqref{eq:excess}--\eqref{eq:hankel} is just the high spectral marginal \(\mu\), for any initial joint state. On the \(m\)-particle high sector with \(m\geq1\), \(H_{0,H}\) multiplies by \(\sum_{a=1}^m k_a\); each level set has Lebesgue measure zero. Therefore the restriction of any normal spectral law to the non-vacuum sectors has no atoms, and the only possible atom of \(R\) is \(r=1\), the high vacuum. Any nonzero weight on non-vacuum sectors gives infinite support and hence both the strict causal-break witness and unbounded exact cut rank. If instead \(\mu=\delta_1\), positivity forces the joint density to be supported on the high-vacuum factor, and \eqref{eq:history} is the ordinary unitary Markov process with the initial low density. Thus a normal physical preparation is operationally memoryless for the stated controls exactly when its high marginal is the vacuum.

\subsection{Matching the auxiliary mean does not match the clock}

Here two preparations have the same \(\rho_L(p)\), the same \(c,p,h_1\), and equal branchwise and total mean auxiliary high energy. For branch zero pick a smooth normalized energy density \(\eta(E)=|\varphi(E)|^2\) supported in a compact interval about \(E_0>\Lambda\). Fix \(\Delta>0\) so that the two translates below have disjoint supports above \(\Lambda\), and choose \(h_1\) supported away from \(\eta\) and both translates. Compare
\[
 \eta_{\rm A}(E)=\eta(E),\qquad
 \eta_{\rm B}(E)=\tfrac12\bigl[\eta(E-\Delta)+\eta(E+\Delta)\bigr].
\]
These are actual one-particle preparations: use \(\varphi\) or the equal norm superposition of its two disjoint translates. The energy moments used here are finite by compact support. The means of \(H_{0,H}\) agree by symmetry of the translates, not by matching a mean of the bounded physical Hamiltonian. For \(r=e^{-bE}\), their moments obey
\[
 m_{{\rm B},n}=m_{{\rm A},n}\cosh(nb\Delta),\qquad
 m_{{\rm A},n}=\int e^{-nbE}\eta(E)\,\dd E .
\]
The moments are conditional and trace-one, since the common initial branch probability can be divided out. The ratio \(m_2/m_1^2\) changes by
\(\cosh(2b\Delta)/\cosh^2(b\Delta)>1\). Hence the two spectral records and Ramsey histories differ. Moreover changing one physical elapsed time by a constant factor would multiply its first two moments by that factor and its square, and could not remove this difference.

In fact the obstruction is a local invariant of the observed complex curve, not just of its time derivatives. Normalize the branch signal \(G_j\) to \(G_j(0)=1\) using the common positive initial branch probability. By \eqref{eq:moments} its imaginary part is nonzero for sufficiently small positive \(t\) and
\begin{equation}
 \lim_{t\downarrow 0}
 \frac{1-\operatorname{Re}G_j(t)}
      {(\operatorname{Im}G_j(t))^2}
 =\frac{\alpha_{j,2}}{2\alpha_{j,1}^2}\frac{m_2}{m_1^2}.
 \label{eq:clock}
\end{equation}
A continuous relabeling of elapsed time that fixes zero and gives positive times arbitrarily close to zero does not change this limit along the signal curve. In particular no common continuous monotone change of clock can reconcile the two short-time Ramsey curves, since \eqref{eq:clock} differs. This uses the same calibrated gap and both quadratures, not a conditional choice of clock after reading the hidden spectral parameter or a change of apparatus.

There is also a pair with both the same auxiliary mean and the same mean of the full physical \(H\). Use four disjoint translates \(\eta(E-j\Delta)\), \(0\leq j\leq3\), all supported above \(\Lambda\) and disjoint from the fixed other branch. Put \(z_*=e^{-b\Delta}\) and define \(c_j\) by
\[
 \sum_{j=0}^3 c_j z^j=P(z)=(z-1)^2(z-z_*).
\]
Take weights \(1/4+\varepsilon c_j\) and \(1/4-\varepsilon c_j\), with \(\varepsilon>0\) small enough for positivity. Each law is again realized by a single smooth packet, the sum of the four disjoint translates of \(\varphi\) with the square roots of the weights. The constraints \(P(1)=P'(1)=0\) give normalization and the same auxiliary-energy mean. The identity \(P(z_*)=0\) gives the same first moment \(m_1\) of \(r\), and hence, with the fixed low branch, the same \(\tr(\rho H)\) by \eqref{eq:physical} and \eqref{eq:assignmentmeasure}. But \(P(z_*^2)\ne0\), so \(m_2\) and the invariant \eqref{eq:clock} still differ. Thus neither of these energy means, alone or together, suffices.

\section{Scope}

The conserved quantity here matters because it multiplies the retained phase in the actual constrained-model Hamiltonian. The complete retained record is neither the high auxiliary mean nor the full joint state: it is a positive low operator for each measurable spectral range. In particular, correlations with a conserved continuous spectrum can survive an arbitrarily explicit causal reset of all retained modes, yet be exactly identifiable by a fixed finite-width gap and varied initial effects. The theorem needs the full stated retained-mode operations and ideal exact time data for that classification. It does not reconstruct environmental coherence outside this spectral algebra, prove a universal spatial boundary sewing law, or extend to arbitrary gravitational polarizations or Hamiltonians whose discarded sector is not conserved in this way.

\begin{thebibliography}{9}
\bibitem{AP}
A.~Ashtekar and M.~Pierri,
\emph{Probing quantum gravity through exactly soluble midi-superspaces I},
J. Math. Phys. \textbf{37}, 6250--6270 (1996),
\href{https://arxiv.org/abs/gr-qc/9606085}{arXiv:gr-qc/9606085}.

\bibitem{BGV}
J.~F.~Barbero G., I.~Garay and E.~J.~S.~Villase\~nor,
\emph{Quantum Einstein--Rosen waves: coherent states and \(n\)-point functions},
Class. Quantum Grav. \textbf{25}, 205013 (2008),
\href{https://arxiv.org/abs/0808.2561}{arXiv:0808.2561}.

\bibitem{PollockProcess}
F.~A.~Pollock, C.~Rodr\'iguez-Rosario, T.~Frauenheim, M.~Paternostro and K.~Modi,
\emph{Non-Markovian quantum processes: Complete framework and efficient characterization},
Phys. Rev. A \textbf{97}, 012127 (2018),
\href{https://doi.org/10.1103/PhysRevA.97.012127}{doi:10.1103/PhysRevA.97.012127}.

\bibitem{PollockMarkov}
F.~A.~Pollock, C.~Rodr\'iguez-Rosario, T.~Frauenheim, M.~Paternostro and K.~Modi,
\emph{Operational Markov condition for quantum processes},
Phys. Rev. Lett. \textbf{120}, 040405 (2018),
\href{https://arxiv.org/abs/1801.09811}{arXiv:1801.09811}.

\bibitem{Goswami}
K.~Goswami, A.~K.~Roy, V.~Srivastava, B.~Perez, C.~Giarmatzi, A.~Gilchrist and F.~Costa,
\emph{Hamiltonian characterization of multi-time processes with classical memory},
New J. Phys. \textbf{27} (2025),
\href{https://doi.org/10.1088/1367-2630/ae1c64}{doi:10.1088/1367-2630/ae1c64};
\href{https://arxiv.org/abs/2412.01998}{arXiv:2412.01998}.

\bibitem{Fanizza}
M.~Fanizza, J.~Lumbreras and A.~Winter,
\emph{Quantum theory in finite dimension cannot explain every general process with finite memory},
Commun. Math. Phys. \textbf{405}, 50 (2024),
\href{https://arxiv.org/abs/2209.11225}{arXiv:2209.11225}.

\bibitem{Pechukas}
P.~Pechukas, \emph{Reduced dynamics need not be completely positive},
Phys. Rev. Lett. \textbf{73}, 1060--1062 (1994),
\href{https://doi.org/10.1103/PhysRevLett.73.1060}{doi:10.1103/PhysRevLett.73.1060}.
\end{thebibliography}
\end{document}
